
\chapter{□□章見出し□□}


□□□□□□□□□□□□□□□□□□□□□□□□□□□□□□□□□□□□□□□□□□□□
□□□□□□□□□□□□□□□□□□□□□□□□□□□□□□□□．

\vspace{0.5\Cvs}
\begin{itemize}
\setlength{\leftskip}{1zw}
 \item □□□□□□□□□□□□□□
 \item □□□□□□□□□□□□□□
 \item □□□□□□□□□□□□□□
\end{itemize}
\vspace{0.5\Cvs}
□□□□□□□□□□□□□□□□□□□□□□□□□□□□．



\section{節見出し□□□□□□□□}
%

□□□□□□□□□□□□□□
□□□□□□□□□□□□□□（{\bf 図\ref{fg:opticalflowimage}}）．□□□□□□□□□□□□□□□□□□□□□□□□□□□□□□□□□□□□□□□□□□
□□□□□□□□□□□□□□□□□□□□□□□□□□□□□□□□□□□□□□□□□□
□□□□□□□□□□□□□□□□□□□□□□□□□□□□□□□□□□□□□□□□□□
□□□□□□□□□□□□□□□□□□□□□□□□□□□□□□□□□□□□□□□□□□
□□□□□□□□□□□□□□□□□□□□□□□□□□□□□□□□□□□□□□□□□□．
%
%%%%%%%%%%%%%%%%%% 図0.1
\begin{figure}[b]
\begin{center}
\includegraphics[scale=1.0]{fig/fig0-1.eps}
 \caption{図版キャプション}
 \label{fg:opticalflowimage}
\end{center}
%\vspace{-2mm}
\end{figure}
%%%%%%%%%%%%%%%%%%
%

\newpage


\subsection{項見出し□□□□□□□□}


\vspace{-6mm}
\begin{progbox}
\begin{verbatim}
void cv::calcOpticalFlowPyrLK( 
        InputArray prevImg, 
        InputArray nextImg,
        InputArray prevPts, 
        InputOutputArray nextPts,
        OutputArray status, 
        OutputArray err
)
\end{verbatim}
\end{progbox}

\noindent{\fontsize{13.5\JQ}{22H}\selectfont\mdgfamily \<\hskip-0.5zw【概要】}

□□□□□□□□□□□□□□□□□□□□□□□□□□□□□□□□□□□□□□□□．
□□□□□□□□□□□□□□□□□□□□□□□□□□□□□□□□□□□□□□□□
□□□□□□□□□□□□□□□□□□□□□□□□□□□□□□□□□□□□□□□□
□□□□□□□□□□□□□□□□□□□□□□□□□□□□□□□□□□□□□□□□
□□□□□□□□
\vspace{0.5\Cvs}

□□□□□□□□□□□□□□□□□□□□□□□□□□□□□□□□□□□□□□□□
□□□□□□□□□□□□□□□□□□□□□□□□□□□□□□□□□□□
□□□□□□□□□□□□□□□□□□□□□□□□□□□□□□□□□□□
□□□□□□□□□□□□□□□□□□□□□□□□□□□□□□□□□□□
□□□□□□□□□□□□□□□□□□□□□□□□□□□□□□□□□□□
□□□□□□□□□□□□□□□□□□□□□□□□□□□□□□□□□□□
□□□□□□□□□□□□□□□□□□□□□□□□□□□□□□□□□□□
\footnote{□□□□□□□□□□□□□□□□□□□□□□□□□□□□□□．}．


%%%%%%%%%%%%%%%%%%%%%%%%%%%%%%%%%%%%%%%%%%%%% プログラムbox
\def\progtitle{sample-prog1(C++)\label{prog1}}
\begin{progboxt}
\lstinputlisting{prog/prog1-4.py}
\end{progboxt}
%%%%%%%%%%%%%%%%%%%%%%%%%%%%%%%%%%%%%%%%%%%%% プログラムbox


\newpage

\shorikekka %%% 処理結果

\vspace{-5mm}

%%%%%%%%%%%%%%%%%% 図
\begin{figure}[h]
\begin{center}
\includegraphics[scale=1.0]{fig/fig0-1.eps}
  \caption{処理結果}
  \label{fg:1}
\end{center}
\vspace{-6mm}
\end{figure}
%%%%%%%%%%%%%%%%%%

\vspace{-2mm}

\def\mamewaku{1}
\begin{mame}
　□□□□□□□□□□□□□□□□□□□□□□□□□□□□□．□□□□□□□□□□□□□□□□□□□□□□□□□□□□□□□
□□□□□□□□□□□□□□□□□□□□□□□□□□□□□□□
□□□□□□□□□□□□□□□□□□□□□□□□□．
\end{mame}
%\vspace{-5mm}


%\newpage
%\hrule width0pt \vspace{-0.7\Cvs}

\begin{renshuu}
 \item[\kmaru{1}] □□□□□□□□□□□□□□□□□□
□□□□□□□□□□□□□□□□□□□□□□□□□□□□□□□□□□□□□□□□．
 \item[\kmaru{2}] □□□□□□□□□□□□□□□□□□□□□□□□□□□□□．   
\end{renshuu}


%{参考文献}
\begin{biblio}
\bibitem{bib:sample1} 
aaaaaaaaaaaaaaaaaaaaaaaaaaaaaaaaaaaaaaaaaaaaaaaaaaaaa,
pp. 363--370, 2003.

\bibitem{bib:sample2}
%無断使用可．書籍の掲載も問題なし．
□□□□□□□□□□□□□□□□□□□□\\
\url{http://index.html}

\end{biblio}

